\section{Main}
\label{sec:introduction}

% Provide the current definition of a phenotype
The National Human Genome Research Institute (NHGRI) defines as phenotype as "an individual’s observable traits, such as height, eye color and blood type" and a trait as "a specific characteristic of an individual." \cite{nhgri_glossary} Traits can be determined by genes, environmental factors or by a combination of both and can be qualitative (such as eye color) or quantitative (such as height or blood pressure). Most genetics and genomics research inevitably attempt to assign genetic associations to observed phenotypes in the population, whether they be complex polygenic traits such as height, or monogenic diseases such as Cystic Fibrosis \cite{riordan1989} or Huntington's Disease \cite{hdcrg1993}. While this has been greatly successful in uncovering numerous genotype to phenptype relationships and yeilding promising theraputics \cite{minikel2024}, we argue the phenotype side of this equation has not be adequately characterized.

% Explain/Illustrate why that definition is a problem
Consider height and weight as two traits of individual, and now consider the concept of Body Mass Index (BMI). The question naturally arises, is BMI a trait? It certianly seems so, as it is a observable characteristic of an individual, but it would seem that the compositionality of the trait should matter, as its result is deterministically derived by height and weight. So can a trait be defined by other traits? In a genome wide association study (GWAS), BMI is certianly treated as a trait \cite{yengo2018}. Further, there is a temporality dimension to traits as an individuals height and weight and therefore BMI will be determined at the precise point of which it is measured. Further, who a trait is measured must also factor in, as a clinical diagnosis and a biobank self-report should probably not be treated as the same trait, although it often is \cite{schoeler2025,hu2025} Therefore, any analysis proposing a link between the genotype and phenotype of an individual will naturally be confounded by at least the time and manner in which observation was made and on the dependency relationship between traits themselves.

% Explain why ontologies are limiting
To date, the most rigorous characterization and organization of the phenome is through the work of disease and experimental ontologies and controlled vocabularies \cite{efo_malone2010,mondo2026,hpo2024,ordo2014} and they have been quite successful in harmonizing biological data \cite{gwascatalog2023} and ensuring researchers are speaking the a unified language \cite{obofoundry2007}. Further, they can represent a degree of compositionality as these ontologies do describe relationships among traits and phenotypes. However, in the case of height, weight, and BMI they do not encode the functional structure of the relationship. For example, the experimental factors ontology (EFO) \cite{efo_malone2010} puts BMI at 7 hops away from body height and only 1 hop away (Supplementary Figure~\ref{fig:efo_bmi_tree}). This structural asymmetry was curated, and thus even any information theortic construction capable of reweighting the symmetry of BMI = weight / height$^2$ is not possible by the construction of the ontology (Supplementary). The only way, such a relation could be rescued structurally is if EFO added an additional relationship linking BMI to height and weight, which would be an additional curation step.

% Brief overview of solution
In the biomedical ontology field and inferred from the the NHGRI definitions of traits and phenotypes, as a field we have defined phenotypes and traits as labels be that EFO and MONDO IDs or free text (eye color = brown). We assert that traits instead should be represented as functions over observables (individuals). In making this conceptual shift, we are then able to form clear compositions of traits and form an algebra to more precisely define phenotypes.

% Contribution and impact of this perspective